\documentclass[11pt,letterpaper]{article}

\usepackage[T1]{fontenc}
\usepackage{lmodern}
\usepackage[margin=0.72in]{geometry}
\usepackage{amsmath,amssymb,mathtools,bm}
\usepackage{array,booktabs,tabularx}
\usepackage{xcolor}
\usepackage{enumitem}
\usepackage{microtype}
\usepackage{fancyhdr}
\usepackage{hyperref}

% The user-local tools/array package is newer than the available LaTeX kernel.
\makeatletter
\def\vcenter@text{\vbox}
\makeatother

\definecolor{navy}{HTML}{17365D}
\definecolor{muted}{HTML}{4A5568}
\hypersetup{
  colorlinks=true,
  linkcolor=navy,
  urlcolor=navy,
  pdftitle={M4 Moment-Hierarchy Closed-Book Workbook},
  pdfauthor={Local Paper-V learning sequence}
}

\pagestyle{fancy}
\fancyhf{}
\fancyhead[L]{\footnotesize $M_4$ moment hierarchy}
\fancyhead[R]{\footnotesize Closed-book workbook}
\fancyfoot[C]{\thepage}
\renewcommand{\headrulewidth}{0.3pt}

\setlength{\parindent}{0pt}
\setlength{\parskip}{4pt}
\setlist[enumerate]{leftmargin=*,itemsep=5pt}
\renewcommand{\arraystretch}{1.18}
\newcolumntype{Y}{>{\raggedright\arraybackslash}X}
\newcommand{\R}{\mathbb R}
\newcommand{\C}{\mathbb C}
\newcommand{\ii}{\mathrm i}
\newcommand{\Tr}{\operatorname{Tr}}
\newcommand{\expect}[1]{\left\langle #1\right\rangle}
\newcommand{\worklines}[1]{%
  \par\vspace{2pt}%
  \foreach \n in {1,...,#1}{\noindent\rule{\textwidth}{0.25pt}\par\vspace{12pt}}%
}
\usepackage{pgffor}

\begin{document}

\begin{center}
{\color{navy}\LARGE\bfseries From the Archive Equations to the $M_4$ Cone}\par
\vspace{3pt}{\Large Closed-Book Reconstruction and Problem Workbook}\par
\vspace{3pt}{\small Post-Reading Mastery Quiz (Exit Quiz) and Transfer-and-Repair Work}
\end{center}
\hrule\vspace{8pt}

Name or session ID: \rule{0.35\textwidth}{0.4pt}\hfill
Date: \rule{0.22\textwidth}{0.4pt}

Start time: \rule{0.18\textwidth}{0.4pt}\hfill
End time: \rule{0.18\textwidth}{0.4pt}\hfill
Sources consulted during attempt: \rule{0.22\textwidth}{0.4pt}

\section*{Administration protocol}

Attempt this workbook after reading
\emph{From the Archive Equations to the $M_4$ Cone}.  Keep the book and the
solutions PDF closed.  Preserve the first answer even when you revise it.
For every item, record elapsed time, confidence, and reasoning.  A correct
choice without a reason is recognition evidence; the derivation and transfer
items test whether the construction can be reproduced independently.

If an item fails, make one second attempt after identifying the earliest
unsupported step.  Open the separate solutions only after that second attempt
or an explicit decision to stop retrieval.  Long derivations may continue on
numbered paper or tablet pages.

\begin{center}
\begin{tabular}{@{}lll@{}}
\toprule
Confidence & Meaning & Record\\
\midrule
0--30\% & recognition or guess & low\\
31--70\% & partial model with uncertainty & medium\\
71--100\% & independently reconstructable & high\\
\bottomrule
\end{tabular}
\end{center}

\section*{Part I --- Governing-map reconstruction}

\textbf{R1. Reconstruct the complete calculation map.}
Begin with \(H(t)\) and the initial quantum state.  End with the time-RMS
observables in the matched result.  Your map must place the archive tuple,
29-coordinate raw chart, 31 hidden moments, $K/P/D$ source, degree-four
completion, retained cone controller, stage retraction, and exact reference in
their correct causal order.  Mark every operation that is offline.

Elapsed: \rule{0.12\textwidth}{0.4pt}\quad
Confidence: \rule{0.12\textwidth}{0.4pt}
\worklines{5}

\textbf{R2. Derive the expectation-value equation.}
Starting only from \(\ii\partial_t\lvert\Psi\rangle=H\lvert\Psi\rangle\),
derive \(\frac{d}{dt}\langle O\rangle\) for a time-independent operator.
State the additional term required when \(O=O(t)\).

Elapsed: \rule{0.12\textwidth}{0.4pt}\quad
Confidence: \rule{0.12\textwidth}{0.4pt}
\worklines{4}

\textbf{R3. Exhibit one hierarchy-generating chain.}
Using the relative-mode Hamiltonian, derive a one-spin expectation whose
derivative requests a spin--phonon moment, then differentiate that mixed moment
until a higher total degree appears.  Label the degree of each moment.

Elapsed: \rule{0.12\textwidth}{0.4pt}\quad
Confidence: \rule{0.12\textwidth}{0.4pt}
\worklines{6}

\textbf{R4. Explain the two coordinate counts.}
Recover the 31-real count of \((\rho,B,N,A,C)\).  Then explain why the APCM
backbone has 29 raw coordinates and why adding its hidden state yields 60.
Identify the exact equality responsible for the 31-to-29 reduction.

Elapsed: \rule{0.12\textwidth}{0.4pt}\quad
Confidence: \rule{0.12\textwidth}{0.4pt}
\worklines{4}

\textbf{R5. Derive moment-matrix positivity.}
Given words \(W_1,\ldots,W_d\), construct \(Q(z)\), define the Gram matrix,
and prove that every physical state gives a positive-semidefinite matrix.
Explain why words through degree two request moments through degree four.

Elapsed: \rule{0.12\textwidth}{0.4pt}\quad
Confidence: \rule{0.12\textwidth}{0.4pt}
\worklines{5}

\textbf{R6. Separate the three cone operations.}
For retained cone slicing, positive fourth-moment completion, and hidden-stage
retraction, state: the optimization variable, the matrix constrained to be
positive semidefinite, and where the operation enters an SSPRK stage.

Elapsed: \rule{0.12\textwidth}{0.4pt}\quad
Confidence: \rule{0.12\textwidth}{0.4pt}
\worklines{5}

\clearpage
\section*{Part II --- Post-Reading Mastery Quiz (Exit Quiz)}
\subsection*{Conceptual judgment}

For each item, circle one answer.  Then write why it defeats the nearest
alternative and record confidence.

\textbf{C1.} A retained moment \(\mu_r\) has total degree \(r\).  Under the
Holstein interaction \(gxZ_s\), its Heisenberg derivative contains a moment of
degree \(r+1\).  Which statement best identifies the resulting modeling
problem?

\begin{enumerate}[label=\Alph*.]
\item The coordinate chart is too small only because Hermitian matrices store
complex entries redundantly.
\item The numerical solver must lower its step until the degree-\(r+1\) moment
vanishes.
\item A finite autonomous model needs a rule that expresses or supplies the
requested degree-\(r+1\) moment from its retained state.
\item The Hamiltonian has left the fixed-particle-number Hilbert space and must
be projected back.
\end{enumerate}
Confidence: \rule{0.12\textwidth}{0.4pt}\quad
Nearest alternative and reason: \rule{0.40\textwidth}{0.4pt}

\textbf{C2.} What does the zero-terminal-cumulant rule establish?

\begin{enumerate}[label=\Alph*.]
\item It gives a deterministic factorized frontier from lower moments, making
the truncated hierarchy executable without certifying positivity or accuracy.
\item It proves that all connected correlations above the retained order
vanish in the exact driven state.
\item It projects the retained state to the nearest point of the $M_4$ cone.
\item It removes all degree-four moments from the degree-three equations.
\end{enumerate}
Confidence: \rule{0.12\textwidth}{0.4pt}\quad
Nearest alternative and reason: \rule{0.40\textwidth}{0.4pt}

\textbf{C3.} Why does the 60-coordinate APCM state contain both archive and
hidden moments?

\begin{enumerate}[label=\Alph*.]
\item The archive coordinates are numerical copies retained only for plotting,
while the hidden moments carry all physical dynamics.
\item The hidden moments encode the exact cutoff-16 wavefunction in compressed
form, while the archive coordinates set the initial state.
\item Both sets are independent copies of the same complete degree-three
hierarchy and should agree componentwise.
\item The archive chart supplies the established low-order backbone, while the
hidden moments supply two-spin and degree-three information that the backbone
does not determine.
\end{enumerate}
Confidence: \rule{0.12\textwidth}{0.4pt}\quad
Nearest alternative and reason: \rule{0.40\textwidth}{0.4pt}

\textbf{C4.} Which statement most precisely defines $M_4$?

\begin{enumerate}[label=\Alph*.]
\item It is the vector of all fourth cumulants, ordered by electronic site.
\item It is the Gram matrix of Pauli--Weyl words through degree two, whose
entries are affine functions of moments through degree four.
\item It is the fourth Runge--Kutta stage matrix used to enforce the retained
joint cone.
\item It is the Hessian of the terminal-completion objective with respect to
the hidden moments.
\end{enumerate}
Confidence: \rule{0.12\textwidth}{0.4pt}\quad
Nearest alternative and reason: \rule{0.40\textwidth}{0.4pt}

\textbf{C5.} Why can positive completion change trajectory accuracy even
though it is written as a representability constraint?

\begin{enumerate}[label=\Alph*.]
\item It replaces the autonomous propagation by exact cutoff-16 contractions
whenever a negative mode appears.
\item It changes only a diagnostic eigenvalue, which indirectly changes the
reported RMS denominator.
\item Its selected degree-four moments enter degree-three velocities, and its
stage retractions can change the accepted hidden lower moments.
\item It forces every retained observable to equal its exact value at each
saved time.
\end{enumerate}
Confidence: \rule{0.12\textwidth}{0.4pt}\quad
Nearest alternative and reason: \rule{0.40\textwidth}{0.4pt}

\textbf{C6.} In the matched factorial ablation, what is the direct task of the
retained joint-Gram controller?

\begin{enumerate}[label=\Alph*.]
\item Choose a nearby velocity that keeps the electronic and retained joint
moment barriers feasible while satisfying declared equalities.
\item Choose the degree-four frontier closest to the zero-cumulant prior.
\item Fit the exact reference derivative to the archive correlation block.
\item Increase the hierarchy order whenever the $C$-block error grows.
\end{enumerate}
Confidence: \rule{0.12\textwidth}{0.4pt}\quad
Nearest alternative and reason: \rule{0.40\textwidth}{0.4pt}

\textbf{C7.} The $M_4$-only and both-cones cells nearly coincide through
\(t t_{\rm hop}=4\), and the $M_4$-only retained Gram minimum remains positive.
What is the strongest licensed conclusion?

\begin{enumerate}[label=\Alph*.]
\item The retained joint-Gram cone is mathematically redundant for every
positive degree-four moment sequence.
\item The retained controller is implemented incorrectly because it produces
almost no visible change.
\item Exact wavefunction propagation is unnecessary whenever $M_4$ remains
positive.
\item Under this matched protocol, positive completion also kept the retained
Gram feasible, so the retained controller added little over this horizon.
\end{enumerate}
Confidence: \rule{0.12\textwidth}{0.4pt}\quad
Nearest alternative and reason: \rule{0.40\textwidth}{0.4pt}

\textbf{C8.} Where does the exact cutoff-16 trajectory enter the matched
comparison?

\begin{enumerate}[label=\Alph*.]
\item Inside every completion solve as the target degree-four frontier.
\item After autonomous rollout, as a common-state observable and error
reference; exact contractions also supply the shared initial preparation.
\item Inside the retained controller whenever the cutting-plane cycle stalls.
\item Only in the figure legend; the table is self-normalized without a
reference.
\end{enumerate}
Confidence: \rule{0.12\textwidth}{0.4pt}\quad
Nearest alternative and reason: \rule{0.40\textwidth}{0.4pt}

\subsection*{Terminology and object roles}

\textbf{T1.} A spectrahedron in the cone-slicing construction is:

\begin{enumerate}[label=\Alph*.]
\item the eigenvector basis of one barrier matrix;
\item the set of all spectra produced by the exact reference trajectory;
\item the reduced-coordinate set for which an affine Hermitian matrix is
positive semidefinite;
\item the span of the degree-four coefficient matrices alone.
\end{enumerate}
Confidence: \rule{0.12\textwidth}{0.4pt}

\textbf{T2.} A hidden-stage retraction is:

\begin{enumerate}[label=\Alph*.]
\item a smallest scaled change to selected hidden lower moments that restores
an extendable positive moment sequence at a provisional integration stage;
\item an interpolation of the stored 31-coordinate trajectory onto the finer
SSPRK grid;
\item a change of the phonon cutoff that removes an unphysical Fock state;
\item a projection of the exact state onto the archive ansatz.
\end{enumerate}
Confidence: \rule{0.12\textwidth}{0.4pt}

\textbf{T3.} An exposed-mode sensitivity \(a_j=v^\dagger\Delta B(e_j)v\)
measures:

\begin{enumerate}[label=\Alph*.]
\item the curvature of the barrier eigenvalue along a complete nonlinear
trajectory;
\item the first-order change of the tested barrier quadratic form per unit
correction coordinate;
\item the overlap between the correction and the exact derivative residual;
\item the probability of occupying the eigenmode \(v\).
\end{enumerate}
Confidence: \rule{0.12\textwidth}{0.4pt}

\textbf{T4.} A conic backward-error tolerance of \(10^{-6}\) means:

\begin{enumerate}[label=\Alph*.]
\item every observable has absolute error below \(10^{-6}\);
\item the correction norm remains below \(10^{-6}\);
\item the exact state has fidelity at least \(1-10^{-6}\);
\item a reported small negative cone eigenvalue within that scale is accepted
as numerical feasibility error, not as a resolved physical excursion.
\end{enumerate}
Confidence: \rule{0.12\textwidth}{0.4pt}

\subsection*{Recognition of mathematical forms}

\textbf{M1.} For time-independent \(O\), which identity follows from the
Schr\"odinger equation?
\[
\begin{aligned}
\mathrm A\quad&\frac{d}{dt}\expect O=-\ii\expect{\{H,O\}},\\
\mathrm B\quad&\frac{d}{dt}\expect O=\ii\expect{[H,O]},\\
\mathrm C\quad&\frac{d}{dt}\expect O=\expect{H}\expect O,\\
\mathrm D\quad&\frac{d}{dt}\expect O=\ii[\expect H,\expect O].
\end{aligned}
\]
Answer: \rule{0.08\textwidth}{0.4pt}\quad Confidence: \rule{0.12\textwidth}{0.4pt}

\textbf{M2.} Which degree matches
\(\expect{X_\uparrow Z_\downarrow x^2p}\)?
\[
\begin{aligned}
\mathrm A\quad&5, &
\mathrm B\quad&4, &
\mathrm C\quad&3, &
\mathrm D\quad&6.
\end{aligned}
\]
Answer: \rule{0.08\textwidth}{0.4pt}\quad Confidence: \rule{0.12\textwidth}{0.4pt}

\textbf{M3.} Which equality proves moment-matrix positivity for
\(Q(z)=\sum_a z_aW_a\)?
\[
\begin{aligned}
\mathrm A\quad&z^\dagger Mz=\expect Q^2,\\
\mathrm B\quad&z^\dagger Mz=\lvert\expect Q\rvert,\\
\mathrm C\quad&z^\dagger Mz=\Tr(Q),\\
\mathrm D\quad&z^\dagger Mz=\expect{Q^\dagger Q}.
\end{aligned}
\]
Answer: \rule{0.08\textwidth}{0.4pt}\quad Confidence: \rule{0.12\textwidth}{0.4pt}

\textbf{M4.} With \(\lambda<0\), which vector solves
\(\min_w\frac12\|w\|_2^2\) subject to
\(a^{\mathsf T}w+\lambda\ge0\)?
\[
\begin{aligned}
\mathrm A\quad&w^\star=\lambda a,\\
\mathrm B\quad&w^\star=-\frac{1}{\lambda}a,\\
\mathrm C\quad&w^\star=-\frac{\lambda}{a^{\mathsf T}a}a,\\
\mathrm D\quad&w^\star=-\lambda(a^{\mathsf T}a)a.
\end{aligned}
\]
Answer: \rule{0.08\textwidth}{0.4pt}\quad Confidence: \rule{0.12\textwidth}{0.4pt}

\textbf{M5.} At reduced iterate \(y^{(k)}\), a normalized mode has
\(B(y^{(k)})v=\lambda v\), \(\lambda<0\), and reduced response \(n\).
Which cut makes this fixed mode nonnegative?
\[
\begin{aligned}
\mathrm A\quad&n^{\mathsf T}y\le n^{\mathsf T}y^{(k)}+\lambda,\\
\mathrm B\quad&n^{\mathsf T}y\ge n^{\mathsf T}y^{(k)}-\lambda,\\
\mathrm C\quad&n^{\mathsf T}y=\lambda,\\
\mathrm D\quad&y^{\mathsf T}n\ge\lambda\,y^{(k)\mathsf T}n.
\end{aligned}
\]
Answer: \rule{0.08\textwidth}{0.4pt}\quad Confidence: \rule{0.12\textwidth}{0.4pt}

\clearpage
\section*{Part III --- Derivation and execution}

\textbf{D1. Cumulant reconstruction.}
Starting from the definition that the third cumulant is the part of
\(\langle ABC\rangle\) not generated by lower partitions, derive the
zero-third-cumulant reconstruction of \(\langle ABC\rangle\).  Then state the
corresponding logical operation at the degree-four frontier of a degree-three
hierarchy.
\worklines{6}

\textbf{D2. Relative oscillator equations.}
Use \([x,p]=\ii\) to derive both lines of
\[
\dot{\langle x\rangle}=\omega\langle p\rangle,
\qquad
\dot{\langle p\rangle}=-\omega\langle x\rangle-g\langle Z_\uparrow+Z_\downarrow\rangle.
\]
Show the two nontrivial commutators.
\worklines{6}

\textbf{D3. One-cut KKT solution.}
Derive the optimizer for
\(\min_w\frac12w^{\mathsf T}w\) subject to
\(a^{\mathsf T}w+\lambda\ge0\), covering both \(\lambda\ge0\) and
\(\lambda<0\).  Interpret the multiplier and the correction direction.
\worklines{7}

\textbf{D4. Equalities before inequalities.}
Given \(Lw=d\), derive the parameterization
\(w=w_{\rm p}+Zy\) from an SVD or pseudoinverse.  State the properties of
\(w_{\rm p}\) and \(Z\) that reduce the Euclidean objective to a constant plus
\(\frac12\|y\|_2^2\).
\worklines{7}

\textbf{D5. Why one slice is insufficient.}
Prove that one exposed eigenvector gives an exact inequality for that fixed
direction.  Then explain, using the Rayleigh quotient, why satisfying that
inequality does not certify the full matrix.  State the termination test that
does certify the implemented tolerance.
\worklines{7}

\textbf{D6. Read the factorial ablation.}
Use the matched table to calculate the improvement factors produced by the
$M_4$ cone alone for: 31-scalar RMS, \(C\)-block RMS, and total-internal-energy
RMS.  Then write one conclusion about mechanism and one conclusion the table
cannot support.
\worklines{6}

\textbf{D7. Insert both cone mechanisms into SSPRK(3,3).}
Write the three SSPRK stages.  Mark where degree-four completion supplies the
velocity, where the retained velocity correction enters, and where hidden
stage retraction occurs.  Explain why checking only the stored end-of-step
state would define a different algorithm.
\worklines{8}

\clearpage
\section*{Part IV --- Transfer-and-Repair Work}

These questions remain unsolved in the companion solutions.  Discuss them in
the later debrief after preserving the first attempt.

\textbf{X1. Change the word basis.}
Suppose the $M_4$ word list omits the mixed words \(X_sx\) and \(X_sp\) while
retaining pure Pauli and pure boson words through degree two.  Which
electron--phonon correlations become invisible to its quadratic-form tests?
Construct an operator combination \(Q(z)\) that the reduced list can no longer
form, and state the likely representability consequence.
\worklines{7}

\textbf{X2. Positivity without accuracy.}
Construct a logically consistent scenario in which \(M_4(t)\succeq0\) for the
entire trajectory while a retained observable has large error against exact
propagation.  Identify which validity axis passes and which one fails.
\worklines{6}

\textbf{X3. Change the correction metric.}
Replace the coordinate-Euclidean norm of \(w\) by the Frobenius norm of its
lifted matrix blocks.  Derive the whitened coordinate transformation that turns
the weighted objective into \(\frac12\|y\|_2^2\).  Explain why the feasible set
is unchanged while the selected correction generally changes.
\worklines{8}

\textbf{X4. Separate the cones.}
Imagine a future run in which positive completion keeps \(M_4\succeq0\), yet
the separately constructed retained joint Gram matrix \(\mathcal G\) becomes
negative.  Give two mathematically distinct explanations: one based on word
content and one based on coordinate/transcription inconsistency.  Design the
first diagnostic for each.
\worklines{8}

\textbf{X5. Faster solver, different trajectory.}
An exact symmetry transformation block-diagonalizes $M_4$, and a faster conic
solver returns a point with essentially the same objective and feasible
eigenvalues.  Its lower-moment retraction differs by \(10^{-4}\), and the
autonomous trajectory later diverges from the serial baseline.  Decide which
tests are required before replacing the baseline and explain why equality of
the formal feasible set is insufficient.
\worklines{8}

\section*{Attempt record}

\begin{tabularx}{\textwidth}{@{}lYYYY@{}}
\toprule
Part & elapsed time & independent items & hinted items & unresolved items\\
\midrule
I: reconstruction &&&&\\[12pt]
II: exit quiz &&&&\\[12pt]
III: derivations &&&&\\[12pt]
IV: transfer &&&&\\[12pt]
\bottomrule
\end{tabularx}

\vspace{12pt}
Earliest recurring unsupported step:
\worklines{2}

Next retrieval date or session:
\rule{0.35\textwidth}{0.4pt}

\end{document}
